% Generated draft: requires mathematical and source review.
\documentclass[11pt]{article}
\usepackage[margin=25mm]{geometry}
\usepackage{fontspec}
\setmainfont{Latin Modern Roman}
\usepackage{amsmath,amssymb,mathtools,unicode-math}
\setmathfont{Latin Modern Math}
\usepackage{enumitem,needspace,xurl,hyperref}
\setlength{\parindent}{0pt}
\setlength{\parskip}{5pt}
\setlength{\emergencystretch}{3em}
\newcommand{\N}{\mathbb{N}}
\newcommand{\Z}{\mathbb{Z}}
\newcommand{\Q}{\mathbb{Q}}
\newcommand{\R}{\mathbb{R}}
\newcommand{\C}{\mathbb{C}}
\newcommand{\F}{\mathbb{F}}
\newcommand{\A}{\mathbb{A}}
\newcommand{\PP}{\mathbb{P}}
\newcommand{\E}{\mathbb{E}}
\newcommand{\eps}{\varepsilon}
\newcommand{\dd}{\,\mathrm{d}}
\begin{document}
\section*{UnsolvedMath ID: 800889}
\textbf{Source problem code:} catalogue entry
\subsection{Definitions and mathematical statement}
The supplied source does not state the hypotheses on \(X\). The formulation below uses the customary interpretation that \(X\) is a compact K\"ahler manifold of complex dimension \(n\); this ambient hypothesis requires confirmation from the source context.

Let \(\omega\) be a K\"ahler form on \(X\), and let \(\alpha\in H^{1,1}(X,\R)\). A K\"ahler current in \(\alpha\) is a closed real \((1,1)\)-current \(S\in\alpha\) satisfying \(S\geq\varepsilon\omega\) for some \(\varepsilon>0\). The class \(\alpha\) is called big if it contains such a current. Its ample locus is
\[
\operatorname{Amp}(\alpha)=\{x\in X:\text{some K\"ahler current in }\alpha\text{ with analytic singularities is smooth near }x\}.
\]

Choose a smooth closed real \((1,1)\)-form \(\theta\in\alpha\). Write \(\operatorname{PSH}(X,\theta)\) for the upper semicontinuous, locally integrable functions \(\varphi\) satisfying \(\theta+dd^c\varphi\geq0\), with \(dd^c=\frac{i}{\pi}\partial\bar\partial\). Put
\[
V_\theta=\left(\sup\{u\in\operatorname{PSH}(X,\theta):u\leq0\}\right)^* ,
\]
where the star denotes upper semicontinuous regularization. For a closed positive current \(T=\theta+dd^c\varphi\), denote its non-pluripolar Monge--Amp\`ere measure by \(\langle T^n\rangle\); this product is constructed by truncating local plurisubharmonic potentials and assigning no mass to pluripolar sets. Define
\[
\operatorname{vol}(\alpha)=\int_X\left\langle(\theta+dd^cV_\theta)^n\right\rangle.
\]
Under the standard \(\mathcal E^1\) interpretation, \(T\) has finite energy if it has full Monge--Amp\`ere mass and finite relative first energy, equivalently
\[
\int_X\langle T^n\rangle=\operatorname{vol}(\alpha),
\qquad
\int_X|\varphi-V_\theta|\,\langle T^n\rangle<\infty.
\]
The integrand is understood almost everywhere for the non-pluripolar measure.

\textbf{Question (Dinew, Guedj, Zeriahi, 2015, Question 20).} Let \(\alpha\) be big and let \(\mu\) be a smooth positive volume form on \(X\). As in the supplied statement, assume that \(T\in\alpha\) is the unique closed positive finite-energy current satisfying
\[
\langle T^n\rangle=\mu.
\]
Must \(T|_{\operatorname{Amp}(\alpha)}\) be represented by a smooth real \((1,1)\)-form?

The assumed existence of a finite-energy solution forces \(\int_X\mu=\operatorname{vol}(\alpha)\). This necessary normalization is not explicitly stated in the supplied excerpt. The compact K\"ahler setting, positivity convention for a volume form, and \(\mathcal E^1\) meaning of finite energy are explicit interpretive qualifications here, not hypotheses verified against the full source.
\subsection{Short English statement}
In the customary compact K\"ahler setting, does the unique finite-energy positive current solving a non-pluripolar Monge--Amp\`ere equation with smooth positive volume data become smooth throughout the ample locus of its big class? The excerpt omits the ambient hypotheses and explicit mass normalization; these require source-context review.
\subsection{Sources}
\begin{itemize}
\item Supplied local UnsolvedMath registry record: permanent ID \(800889\), exact source problem code \texttt{catalogue entry}, title \emph{Open Problems in Pluripotential Theory --- Question 20}. The record attributes the question to Dinew, Guedj, Zeriahi, with proposed year 2015.
\item Dinew, Guedj, Zeriahi, \emph{Open Problems in Pluripotential Theory} (2015), Question 20. Supplied source URL: \url{https://arxiv.org/abs/1511.00705}. The supplied provenance records access on 2026-10-16 and extraction by \texttt{arxiv\_latex}; these are supplied provenance statements, not independent verification here.
\item The registry labels the problem \texttt{open} and records an update dated 2026-10-16. This is a supplied status label, not an independently verified assessment of the literature.
\end{itemize}
\end{document}
